

\section{Metrics on Probability Measures: From Historical Origins to Statistical Applications}\label{section:general_intro_metrics_on_probabilities}

Historically, several distinct motivations led to the development of distances between probability measures. The first motivation is purely mathematical, with formal foundations laid following the introduction of metric spaces by Felix Hausdorff and Maurice Fr\'echet at the beginning of the 20th century. This was followed by work on the space of signed measures that led to the development of the total variation distance. In parallel, a metric for distribution functions on the real line was introduced in \cite{levy1925calcul}. In addition, in a note written for the book \cite{frechet1950recherches}, L\'evy introduced several notions of distance between random variables and between their induced laws, one of which corresponds to what is now recognized as the Wasserstein distance, with emphasis on distances between random variables and their probability laws.

Another motivation for studying distances between probability measures arose in information theory. In \cite{kullback1951information}, the notion of information for discrimination between two distributions was introduced. The motivation was to quantify how much information observations provide for distinguishing between two hypotheses about the underlying distribution. The information-theoretic point of view was followed by a broader development of divergence measures, such as the chi-square divergence, \(f\)-divergences, and the squared Hellinger distance.

One of the most widely used distances on the space of probability measures is the Wasserstein distance. Although its motivation was not the study of metrics on the space of probability measures, its origins can be traced back to the development of optimal transport theory. In \cite{monge1781memoire}, the practical question of how to move a pile of soil to fill another hole while minimizing the total transportation cost was posed. The problem was formulated in terms of transport maps, which specify where each infinitesimal unit of mass is sent. Around 160 years later, a relaxation of this formulation was introduced in \cite{kantorovich1942translocation}, allowing transport plans rather than deterministic maps. This relaxation led to what is now known as the Monge--Kantorovich problem:
\begin{equation}\label{equation:kantorovich_formulation}
\min_{\gamma \in \Gamma(\mu, \nu)}
\int_{\X\times \Y} c(x,y)\,\mathrm{d}\gamma(x,y),
\end{equation}
where \(\Gamma(\mu,\nu)\) denotes the set of all couplings of \(\mu\) and \(\nu\), that is, probability measures on \(\X\times\Y\) with marginals \(\mu\) and \(\nu\). This formulation was later complemented by the duality results of \cite{kantorovich1958space}.

The Monge--Kantorovich problem includes the Wasserstein distance as a special case. Indeed, when \(\X=\Y\) is equipped with a metric \(d\), and the cost is chosen as \(c(x,y)=d(x,y)^p\) for \(p\geq 1\), the corresponding optimal value defines the \(p\)-Wasserstein distance,
\[
W_p(\mu,\nu)
=
\left(
\inf_{\gamma\in\Gamma(\mu,\nu)}
\int_{\X\times\Y} d(x,y)^p\,\mathrm d\gamma(x,y)
\right)^{1/p}.
\]
However, the terminology itself has a separate historical origin. The name ``Wasserstein distance'' first appeared in \cite{dobrushin1970prescribing}, in the context of studying the existence and uniqueness of random fields, as a reference to earlier work by Leonid Vaserstein. In \cite{vaserstein1969markov}, Markov processes on denumerable products of spaces were studied, motivated by the analysis of large systems of automata, and that distance was used to study the convergence of Markov processes.

Wasserstein distances have become important tools in statistics and machine learning because they provide a geometrically meaningful discrepancy between probability measures. In statistics, they are used in the theory of point estimation, where parameters are estimated by minimizing the Wasserstein distance between the empirical distribution and a model distribution \cite{bassetti2006wasserstein,bernton2019parameter}; the convergence of empirical estimators based on the Wasserstein distance is studied in \cite{bobkov2019one,panaretos2019statistical}; in hypothesis testing, where Wasserstein distances are used to measure departures from a null hypothesis \cite{ramdas2017wasserstein,wang2021two,hu2025two}; and in goodness-of-fit testing \cite{hallin2021multivariate}. In computational statistics, Wasserstein distances are used to study convergence rates of MCMC algorithms \cite{qin2022wasserstein,lavenant2024convergence} and to design likelihood-free inference methods, such as Wasserstein approximate Bayesian computation, by comparing observed and simulated empirical distributions directly \cite{bernton2019abc}. In machine learning, they play a central role in generative modeling, especially in Wasserstein generative adversarial networks \cite{arjovsky2017wasserstein}, and have also been used in domain adaptation \cite{courty2016optimal}.

